\documentclass[10pt,a4paper]{amsart}
\usepackage[margin=19mm]{geometry}
\usepackage{amsmath,amssymb,mathtools}
\usepackage[scr=rsfs,cal=euler]{mathalfa}
\usepackage{xfrac}
\usepackage[unicode,hidelinks,bookmarks=false]{hyperref}
\newcommand{\ie}{i.e.\ }
\newcommand{\cf}{cf.\ }
\newcommand{\resp}{respectively\ }
\newcommand{\Subsection}{\subsection}
\providecommand{\nobreakdash}{\nobreak\hbox{-}}
\setlength{\emergencystretch}{2em}
\multlinegap0pt
\usepackage{amssymb}
\usepackage{mathtools}
\newtheorem{lemma}{Lemma}
\newtheorem{proposition}{Proposition}
\def\XXint#1#2#3{{\setbox0=\hbox{$#1{#2#3}{\int}$ }
		\vcenter{\hbox{$#2#3$ }}\kern-.6\wd0}}
\def\Xint#1{\mathchoice
	{\XXint\displaystyle\textstyle{#1}}
	{\XXint\textstyle\scriptstyle{#1}}
	{\XXint\scriptstyle\scriptscriptstyle{#1}}
	{\XXint\scriptscriptstyle\scriptscriptstyle{#1}}
	\!\int}
\def\dashint{\Xint-}
\title{Weak Harnack inequality and Cartan property for nonlocal $W^{s,1}$-minimizers}
\author{Panu Lahti, Yuxin Li, Khanh Nguyen}
\date{Source: arXiv:2603.21121v1 (CC BY 4.0)}
\begin{document}
\maketitle

\noindent\emph{Source and licence.} Weak Harnack inequality and Cartan property for nonlocal $W^{s,1}$-minimizers. Version arXiv:2603.21121v1 (CC BY 4.0), distributed by arXiv under the Creative Commons Attribution 4.0 International licence, http://creativecommons.org/licenses/by/4.0/, as recorded in the arXiv licence metadata for this exact version and shown on the abstract page https://arxiv.org/abs/2603.21121v1. Full licence notice: \texttt{licenses/arxiv-2603.21121-CC-BY-4.0.txt}.\par\smallskip

\noindent\textit{Editorial scope.} Counted unit: the article's proposition lem2.2-13oct (source0.tex lines 1213--1238). The article's complete original proof of it is delivered with it (source0.tex lines 1240--1435, one \texttt{proof} environment of 196 lines). No mathematics is written, repaired or re-derived: the statement and the proof are the article's own lines, in the article's own notation, and no retained line is substituted or re-flowed.  Retained uncounted as the source-local context the proof needs: lines 391--395; lines 704--711; lines 1178--1191.  Nothing was omitted from any retained span, so the package carries no omission record.  No retained line contains a citation, so the wrapper carries no bibliography and no bibliography entry.  No retained line contains a figure environment, an image-inclusion command or drawing code, so no image is delivered and the package declares no asset dependency.  Editorial formatting only, in addition: an AMS \texttt{amsart} preamble that restates the article's own declarations and macros used by the retained lines, namely the environments \texttt{lemma}, \texttt{proposition} on separate counters, together with the standard packages those lines need.  The retained environments print in this document as Lemma 1, Proposition 1 in order of appearance, whereas the article prints the same items with the numbering of its own full document.  The complete native source of the article as distributed in the arXiv e-print for this exact version is bundled unchanged as \texttt{sources/196917/source0.tex}.
	\begin{align}\label{eq2.1-11dec}
		\dashint_B |u(x)|\,d\mu(x)
		&\leq  |u(y)|+\dashint_{B} {|u(x)-u(y)|}\,d\mu(x)\\
		&\leq |u(y)|+\, \frac{\mu(4B) }{\mu(B)}(2r)^s\int_{B} \frac{|u(x)-u(y)|}{\mu(B_{x,y})d(x,y)^s}\,d\mu(x)<\infty,  \notag
	\end{align}

		\begin{align}
			&\quad \int_{X}  \frac{  |\eta(y)-\eta(x)|}{\mu(B_{x,y}) d(x,y)^s}  \,d\mu(x)\notag \\
			&= \ \sum_{j=-\infty}^0 \int_{A_j} \frac{  |\eta(y)-\eta(x)|}{\mu(B_{x,y}) d(x,y)^s}  \,d\mu(x) + \sum_{j=1}^\infty \int_{A_j} \frac{  |\eta(y)-\eta(x)|}{\mu(B_{x,y}) d(x,y)^s}  \,d\mu(x) \notag \\
			&\leq   \sum_{j=-\infty}^0 \int_{A_j} \frac{ \frac{2}{r_1-r}d(x,y) }{\mu(B_j) d(x,y)^s}  \,d\mu(x) +  \sum_{j=1}^\infty \int_{A_j} \frac{2}{\mu(B_j) d(x,y)^s}  \,d\mu(x) \notag \\
			&\simeq    \frac{1}{r_1-r}  \sum_{j=-\infty}^0 (2^{j}(r_1-r))^{1-s} +  \sum_{j=0}^\infty (2^j(r_1-r))^{-s}\notag \\
			&\simeq    \frac{1}{(r_1-r)^s} \left( \frac{1}{1-2^{-(1-s)}} + \frac{1}{1-2^{-s}} \right)\notag \\
			&\lesssim \frac{1}{s(1-s)}\frac{1}{(r_1-r)^s}, \label{eq2.6-29oct}
		\end{align}

	\begin{lemma}\label{lem1.3-1oct}
		If $f$ is an $s$-subminimizer 
		{in a given nonempty open set $\Omega$,}
		where $0<s<1$, then
		for every nonpositive  function  $\eta\in W^{s,1}_0(\Omega)$,
		{
			\begin{align*}
				&0\geq  \left( \int_\Omega\int_\Omega +2\int_\Omega\int_{\Omega^c} \right) \frac{ \frac{u(x)-u(y)}{|u(x)-u(y)|} (\eta(y)-\eta(x)) }{\mu(B_{x,y}) d(x,y)^s} 
				{\bf 1}_{\{ u(x)\neq u(y)\}} \,d\mu(x)\,d\mu(y)\\
				&\quad  - \left( \int_\Omega\int_\Omega +2\int_\Omega\int_{\Omega^c} \right) \frac{  |\eta(x)-\eta(y)| }{\mu(B_{x,y}) d(x,y)^s} {\bf 1}_{\{ u(x)= u(y)\}} \,d\mu(x)\,d\mu(y).
			\end{align*}
		}
		
	\end{lemma}

	\begin{proposition}\label{lem2.2-13oct}
		
		
		Let $\Omega$ be a nonempty open set.
		Let $k\in\mathbb{R}$, $x_0\in \Omega$, $0<\rho<R$ with $B(x_0,R)\subset \Omega$, and set $r:=\frac{R+\rho}{2}$.
		Assume that either
		\begin{enumerate}
			\item \label{item-a}
			$u$ is an $s$-subminimizer 
			{in $\Omega$; or }
			\item \label{item-b}
			$u$ is a $\mathcal{K}_{\psi,0}(\Omega)$-solution where $\psi:\Omega\to[-\infty,\infty)$ is a function such that $\psi\leq k$
			a.e.\ on $B(x_0,R)$.
		\end{enumerate}
		Let $\phi$ be a $\frac{1}{r-\rho}$-Lipschitz function with $0\leq \phi\leq 1$, $\phi\equiv 1 $ on $B(x_0,\rho)$ and $\phi\equiv  0$ on $ X\setminus B(x_0,r)$.
		Then we have the following Caccioppoli--De~Giorgi type inequality:  
		there exists a constant $C>0$, depending only on the doubling constant, such that
		\begin{equation}\label{eq3.2-15mar}
			\|(u-k)_+\phi\|_{\dot W^{s,1}(X)}
			\leq 
			\frac{C}{s(1-s)}
			\frac{1}{(R-\rho)^s}
			\int_{B(x_0,R)} (u-k)_+\, d\mu.
		\end{equation}
		
	\end{proposition}

	\begin{proof}
		Denote $B:=B(x_0,R)$. 
		Let $\eta:=-\phi (u-k)_+$.  We first prove that if  either  \ref{item-a} or \ref{item-b} holds, then 
		{
			\begin{align}\notag
				0\geq&\, \left( \int_\Omega\int_\Omega +2\int_\Omega\int_{\Omega^c} \right) \frac{\frac{u(x)-u(y)}{|u(x)-u(y)|} (\eta(y)-\eta(x)) }{\mu(B_{x,y}) d(x,y)^s} {\bf 1}_{\{u(x)\neq u(y) \}} \,d\mu(x)\,d\mu(y)\\
				&\, - \left( \int_\Omega\int_\Omega +2\int_\Omega\int_{\Omega^c} \right) \frac{ |\eta(x)-\eta(y)| }{\mu(B_{x,y}) d(x,y)^s} {\bf 1}_{\{u(x)=u(y)\}}\,d\mu(x)\,d\mu(y) \label{eq3.2-20nov}.
			\end{align}
		}
		By Lemma \ref{lem1.3-1oct}, if $u$ satisfies \ref{item-a}, then  \eqref{eq3.2-20nov}  holds.
		We now consider the case \ref{item-b}. Let $u$ be a $\mathcal{K}_{\psi,0}(\Omega)$-solution,
		and $\psi\leq k$ a.e. on $B(x_0,R)$. Hence $u\in \dot W^{s,1}(X)$,
		$u\geq \psi$ a.e. on $\Omega$, and $u=0$ a.e. on $X\setminus\Omega$.
		Let $0<t<\frac{1}{2}$. Set $v:=u+t\eta=u-t\phi(u-k)_+$. Then for a.e. $x\in\Omega$,
		\begin{align*}
			v(x)& =  \begin{cases}
				(1-t\phi(x))u(x)+t\phi(x) k& {\rm \ \ if\ \ }u(x)>k {\rm \ and\ }x\in B,\\
				u(x)& {\rm \ \ otherwise\ \ }
			\end{cases}\\
			&\geq   \psi(x)
		\end{align*}
		since $u\geq \psi$ a.e. on $\Omega$, $k\geq \psi$ a.e. on $B$, and $0\leq t\phi\leq 1$.
		Moreover, we have $v\in \dot W^{s,1}(X)$ because 
		\[
		\|v\|_{\dot W^{s,1}(X)}
		\leq 
		{ \|u\|_{\dot W^{s,1}(X)} }
		+\|t\phi(u-k)_+\|_{\dot W^{s,1}(X)}
		\leq 2 \|u\|_{\dot W^{s,1}(X)}<\infty.
		\]
		Notice that   $v=u=0$ a.e. on $X\setminus\Omega$. 
		Combining this with the two above properties, we obtain that
		$v\in \mathcal K_{\psi,0}(\Omega)$. Hence since $u$ is a $\mathcal{K}_{\psi,0}(\Omega)$-solution,
		we have $\|u\|_{\dot W^{s,1}(X)}\leq \|v\|_{\dot W^{s,1}(X)}=\|u+t\eta\|_{\dot W^{s,1}(X)}$.
		Since $v=u=0$ a.e. on $X\setminus\Omega$, it follows that 
		{ $0\geq F(u,u+t\eta)$ for  all $0<t<\frac{1}{2}$ because 
			\begin{align*}
				0
				&\geq  \int_X\int_X \frac{|u(x)-u(y)|-|v(x)-v(y)|}{\mu(B_{x,y})d(x,y)^s}\,d\mu(x)\,d\mu(y)\\
				&=\left(\int_{\Omega}\int_{\Omega}+2\int_{\Omega}\int_{\Omega^c} \right)\frac{|u(x)-u(y)|-|v(x)-v(y)|}{\mu(B_{x,y})d(x,y)^s}\,d\mu(x)\,d\mu(y)\\
				& = F(u,u+t\eta)
			\end{align*}
			for all $0<t<\frac{1}{2}$. 
		}
		Using this, the proof of Lemma \ref{lem1.3-1oct} gives the inequality \eqref{eq3.2-20nov}.
		
		
		
		
		Next, by the fact that $\eta\equiv 0$ on $B^c$, we have
		{
			\begin{align*}
				&- \left( \int_\Omega\int_\Omega +2\int_\Omega\int_{\Omega^c} \right) \frac{ |\eta(x)-\eta(y)| }{\mu(B_{x,y}) d(x,y)^s} {\bf 1}_{\{u(x)=u(y)\}}\,d\mu(x)\,d\mu(y)\\
				&=-  \left( \int_B\int_{B} +2\int_{B}\int_{B^c} \right) \frac{ |\eta(x)-\eta(y)| }{\mu(B_{x,y}) d(x,y)^s} {\bf 1}_{\{u(x)=u(y)\}}\,d\mu(x)\,d\mu(y)\\
				&= -2\int_B\int_{B}  \frac{ \eta(y)-\eta(x) }{\mu(B_{x,y}) d(x,y)^s} {\bf 1}_{\{u(x)=u(y)\}}{\bf 1}_{\{ \eta(x)\leq \eta(y)\}}\,d\mu(x)\,d\mu(y)\notag \\
				&\quad   -2\int_{B} \int_{B^c} \frac{ \eta(y) }{\mu(B_{x,y}) d(x,y)^s} {\bf 1}_{\{u(x)=u(y)\}}\,d\mu(x)\,d\mu(y)
			\end{align*}
			and
			\begin{align*}
				&\left( \int_\Omega\int_\Omega+2\int_{\Omega}\int_{\Omega^c}\right) \frac{\frac{u(x)-u(y)}{|u(x)-u(y)|} (\eta(y)-\eta(x)) }{\mu(B_{x,y}) d(x,y)^s} {\bf 1}_{\{ u(x)\neq u(y)\}}\,d\mu(x)\,d\mu(y)\notag \\
				&= \left( \int_B\int_B+2\int_{B}\int_{B^c}\right) \frac{\frac{u(x)-u(y)}{|u(x)-u(y)|} (\eta(y)-\eta(x)) }{\mu(B_{x,y}) d(x,y)^s} {\bf 1}_{\{ u(x)\neq u(y)\}}\,d\mu(x)\,d\mu(y)\notag \\
				&= \int_B\int_B\frac{\frac{u(x)-u(y)}{|u(x)-u(y)|} (\eta(y)-\eta(x)) }{\mu(B_{x,y}) d(x,y)^s} {\bf 1}_{\{ u(x)\neq u(y)\}}\,d\mu(x)\,d\mu(y)  \notag  \\
				&\quad +2\int_{B}\int_{B^c} \frac{\frac{u(x)-u(y)}{|u(x)-u(y)|} \eta(y) }{\mu(B_{x,y}) d(x,y)^s} {\bf 1}_{\{ u(x)\neq u(y)\}}\,d\mu(x)\,d\mu(y).
			\end{align*}
			By \eqref{eq3.2-20nov}, it follows that
			\begin{align}\label{eq1-2oct}
				0\geq T_{\rm L} + T_{\rm NL}
			\end{align}
			where 
			\begin{align*}
				T_{\rm L} &:=  \int_B\int_B\frac{\frac{u(x)-u(y)}{|u(x)-u(y)|} (\eta(y)-\eta(x)) }{\mu(B_{x,y}) d(x,y)^s} {\bf 1}_{\{ u(x)\neq u(y)\}}\,d\mu(x)\,d\mu(y) \\
				& \quad -2\int_B\int_{B}  \frac{ \eta(y)-\eta(x) }{\mu(B_{x,y}) d(x,y)^s} {\bf 1}_{\{u(x)=u(y)\}}{\bf 1}_{\{ \eta(x)\leq \eta(y)\}}\,d\mu(x)\,d\mu(y)
			\end{align*}
			and
			\begin{align*}
				T_{\rm NL} &:= 2\int_{B}\int_{B^c} \frac{\frac{u(x)-u(y)}{|u(x)-u(y)|} \eta(y) }{\mu(B_{x,y}) d(x,y)^s} {\bf 1}_{\{ u(x)\neq u(y)\}}\,d\mu(x)\,d\mu(y)\\
				& \quad -2\int_{B} \int_{B^c} \frac{ \eta(y) }{\mu(B_{x,y}) d(x,y)^s} {\bf 1}_{\{u(x)=u(y)\}}\,d\mu(x)\,d\mu(y).
			\end{align*}
		}
		We need to estimate the ``local'' term $T_{\mathrm{L}}$ and the ``nonlocal''
		term $T_{\mathrm{NL}}$.  
		{We  estimate the first term of $T_{\rm L}$, denoted $T_{\rm L,1}$, that} 
		\begin{align*}
			T_{\rm L, 1} & := \int_B\int_B\frac{\frac{u(x)-u(y)}{|u(x)-u(y)|} (\eta(y)-\eta(x)) }{\mu(B_{x,y}) d(x,y)^s} {\bf 1}_{\{ u(x)\neq  u(y)\}} \,d\mu(x)\,d\mu(y) \\
			&=\int_B\int_B\frac{\frac{u(x)-u(y)}{|u(x)-u(y)|} (\eta(y)-\eta(x)) }{\mu(B_{x,y}) d(x,y)^s} {\bf 1}_{\{ u(x)> u(y)\}} \,d\mu(x)\,d\mu(y)  \notag \\
			&\quad +\int_B\int_B\frac{\frac{u(x)-u(y)}{|u(x)-u(y)|} (\eta(y)-\eta(x)) }{\mu(B_{x,y}) d(x,y)^s} {\bf 1}_{\{ u(x)< u(y)\}} \,d\mu(x)\,d\mu(y) \notag  \\
			&=  \int_B\int_B\frac{\eta(y)-\eta(x) }{\mu(B_{x,y}) d(x,y)^s} {\bf 1}_{\{ u(x)> u(y)\}} \,d\mu(x)\,d\mu(y)  \notag \\
			&\quad +\int_B\int_B\frac{ \eta(x)-\eta(y) }{\mu(B_{x,y}) d(x,y)^s} {\bf 1}_{\{ u(x)< u(y)\}} \,d\mu(x)\,d\mu(y)\notag \\
			&= 2  \int_B\int_B \frac{\eta(x)-\eta(y) }{\mu(B_{x,y}) d(x,y)^s} {\bf 1}_{\{ u(x)< u(y)\}} \,d\mu(x)\,d\mu(y)\notag 
		\end{align*}
		by a change of variables. From this, we have that
		\begin{align}
			T_{\rm L,1}
			&= 2  \int_B\int_B \frac{\eta(x)-\eta(y) }{\mu(B_{x,y}) d(x,y)^s} {\bf 1}_{\{ u(x)< u(y)\}} {\bf 1}_{\{\eta(x)\geq \eta(y)\}} \,d\mu(x)\,d\mu(y)\notag \\
			&\quad + 2  \int_B\int_B \frac{\eta(x)-\eta(y) }{\mu(B_{x,y}) d(x,y)^s} {\bf 1}_{\{ u(x)< u(y)\}} {\bf 1}_{\{\eta(x)\leq \eta(y)\}} \,d\mu(x)\,d\mu(y)\notag \\
			&=  2\int_B\int_B \frac{\eta(x)-\eta(y) }{\mu(B_{x,y}) d(x,y)^s} {\bf 1}_{\{ \eta(x)\geq \eta(y)\}} \,d\mu(x)\,d\mu(y)\notag \\ 
			&\quad - 2  \int_B\int_B \frac{\eta(x)-\eta(y) }{\mu(B_{x,y}) d(x,y)^s} {\bf 1}_{\{ u(x)= u(y)\}} {\bf 1}_{\{\eta(x)\geq \eta(y)\}} \,d\mu(x)\,d\mu(y) \notag \\
			&\quad - 2  \int_B\int_B \frac{\eta(x)-\eta(y) }{\mu(B_{x,y}) d(x,y)^s} {\bf 1}_{\{ u(x)> u(y)\}} {\bf 1}_{\{\eta(x)\geq \eta(y)\}} \,d\mu(x)\,d\mu(y) \notag \\
			&\quad + 2  \int_B\int_B \frac{\eta(x)-\eta(y) }{\mu(B_{x,y}) d(x,y)^s} {\bf 1}_{\{ u(x)< u(y)\}} {\bf 1}_{\{\eta(x)\leq \eta(y)\}} \,d\mu(x)\,d\mu(y) \notag \\
			&=   \int_B\int_B \frac{|\eta(x)-\eta(y)| }{\mu(B_{x,y}) d(x,y)^s}  \,d\mu(x)\,d\mu(y)\notag \\ 
			&\quad - 2  \int_B\int_B \frac{\eta(x)-\eta(y) }{\mu(B_{x,y}) d(x,y)^s} {\bf 1}_{\{ u(x)= u(y)\}} {\bf 1}_{\{\eta(x)\geq \eta(y)\}} \,d\mu(x)\,d\mu(y) \notag \\
			&\quad + 4 \int_B\int_B \frac{\eta(x)-\eta(y) }{\mu(B_{x,y}) d(x,y)^s} {\bf 1}_{\{ u(x)< u(y)\}} {\bf 1}_{\{\eta(x)\leq \eta(y)\}} \,d\mu(x)\,d\mu(y) .\notag
		\end{align}
		{
			Hence
			\begin{align}
				\notag T_{\rm L}&= T_{\rm L,1}   -2\int_B\int_{B}  \frac{ \eta(y)-\eta(x) }{\mu(B_{x,y}) d(x,y)^s} {\bf 1}_{\{u(x)=u(y)\}}{\bf 1}_{\{ \eta(x)\leq \eta(y)\}}\,d\mu(x)\,d\mu(y)\\
				& \geq  \int_B\int_B \frac{|\eta(x)-\eta(y)| }{\mu(B_{x,y}) d(x,y)^s}  \,d\mu(x)\,d\mu(y)\notag\\
				\label{eq2-2oct} & \quad - 8 \int_B\int_B \frac{\eta(y)-\eta(x) }{\mu(B_{x,y}) d(x,y)^s} {\bf 1}_{\{ u(x)\leq  u(y)\}} {\bf 1}_{\{\eta(x)\leq \eta(y)\}} \,d\mu(x)\,d\mu(y).
			\end{align}
		}
		Next, for the non-local term $T_{\rm NL}$,
		we estimate 
		\begin{align*}
			T_{\rm NL}&\geq  -2 \int_{B} \eta(y) \left(\int_{B^c} \frac{1}{\mu(B_{x,y}) d(x,y)^s}\,d\mu(x) \right) d\mu(y)\\
			&=  -2 \int_{B} (u-k)_+(y)\phi(y) \left(\int_{B^c} \frac{1}{\mu(B_{x,y}) d(x,y)^s}\,d\mu(x) \right) d\mu(y).
		\end{align*}
		Setting $r_y=R-d(x_0,y)$ for each $y\in B$, we have that  $B(y,r_y)\subset B$
		for every $y\in B$. It follows
		from the doubling property that  for every $y\in B$,
		\begin{align}\label{eq2-3Lemma2}
			\int_{B^c} \frac{1}{\mu(B_{x,y}) d(x,y)^s}\,d\mu(x)
			&\leq \int_{X\setminus B(y,r_y)} \frac{d\mu(x)}{\mu(B_{x,y}) d(x,y)^s}\notag\\
			&=  \sum_{j=0}^\infty \int_{B(y,\,2^{j+1}r_y)\setminus B(y,\, 2^jr_y)} \frac{d\mu(x)}{\mu(B_{x,y}) d(x,y)^s}\notag\\
			&\simeq   \sum_{j=0}^\infty \frac{1}{(2^j r_y)^s}= \frac{1}{1-2^{-s}}\frac{1}{r_y^s}\notag\\
			& \lesssim \frac{1}{s}\frac{1}{(R-d(x_0,y))^s}.
		\end{align}
		From the last two estimates, we get
		\begin{align}
			T_{\rm NL}
			&\gtrsim  - \frac{1}{s}\int_{B} \frac{(u-k)_+(y)\phi(y)}{(R-d(x_0,y))^s}\,d\mu(y) \notag\\ 
			&=  - \frac{1}{s}\int_{B(x_0,r)} \frac{(u-k)_+(y)\phi(y)}{(R-d(x_0,y))^s}\,d\mu(y)
			\quad\textrm{since $\phi|_{X\setminus B(x_0,r)}\equiv 0$} \notag\\ 
			& \geq - \frac{2^s}{s}\int_{B(x_0,r)} \frac{(u-k)_+(y)\phi(y)}{(R-\rho)^s}\,d\mu(y), \notag
		\end{align}
		because for every $y\in B(x_0,r)$, we have from  $r=\frac{R+\rho}{2}$ that
		\[
		R-d(x_0,y)\geq R-r = \frac{R-\rho}{2}.
		\]
		Combining this with \eqref{eq1-2oct}-\eqref{eq2-2oct}, we get
		{
			\begin{align*}
				0\gtrsim &  \int_B\int_B \frac{|\eta(x)-\eta(y)| }{\mu(B_{x,y}) d(x,y)^s}  \,d\mu(x)\,d\mu(y)\\
				&-   \int_B\int_B \frac{\eta(y)-\eta(x) }{\mu(B_{x,y}) d(x,y)^s} {\bf 1}_{\{ u(x)\leq u(y)\}} {\bf 1}_{\{\eta(x)\leq \eta(y)\}} \,d\mu(x)\,d\mu(y)\\
				& - \frac{1}{s} \int_{B} \frac{(u-k)_+(x)\phi(x)}{(R-\rho)^s} \,d\mu(x).
			\end{align*}
		}
		We have
		\begin{align*}
			&  \int_B\int_B \frac{\eta(y)-\eta(x) } {\mu(B_{x,y}) d(x,y)^s} {\bf 1}_{\{ u(x)\leq u(y)\}} {\bf 1}_{\{\eta(x)\leq \eta(y)\}} \,d\mu(x)\,d\mu(y)\\
			&=  \int_B\int_B \frac{(u-k)_+(y) (\phi(x)-\phi(y))}{\mu(B_{x,y}) d(x,y)^s} {\bf 1}_{\{ u(x)\leq u(y)\}} {\bf 1}_{\{\eta(x)\leq \eta(y)\}} \,d\mu(x)\,d\mu(y)\\
			&\quad +  \int_B\int_B \frac{((u-k)_+(x) -(u-k)_+(y))\phi(x)}{\mu(B_{x,y}) d(x,y)^s} {\bf 1}_{\{ u(x)\leq u(y)\}} {\bf 1}_{\{\eta(x)\leq \eta(y)\}} \,d\mu(x)\,d\mu(y)\\
			&\leq   \int_B\int_B \frac{\max\{(u-k)_+(x), (u-k)_+(y) \} |\phi(y)-\phi(x)|}{\mu(B_{x,y}) d(x,y)^s}  \,d\mu(x)\,d\mu(y),
		\end{align*}
		where the last line is obtained because the second term in the third line is negative. Combining the last two estimates, we conclude that		
		{
			\begin{align}
				0\gtrsim&  \int_B\int_B \frac{|\eta(x)-\eta(y)| }{\mu(B_{x,y}) d(x,y)^s}  \,d\mu(x)\,d\mu(y) \notag \\ \notag
				&-   \int_B\int_B \frac{\max\{(u-k)_+(x), (u-k)_+(y) \} |\phi(y)-\phi(x)|}{\mu(B_{x,y}) d(x,y)^s}  \,d\mu(x)\,d\mu(y)
				\\ & - \frac{1}{s} \int_{B} \frac{ (u-k)_+(x)\phi(x) }{(R-\rho)^s}\,d\mu(x).\label{eq2.3-13oct}
			\end{align}
		}
		Notice that $\phi$ is $\frac{1}{r-\rho}$-Lipschitz, and so $\frac{2}{R-\rho}$-Lipschitz because $r=\frac{R+\rho}{2}$. We have from \eqref{eq2.6-29oct} that
		\[
		\int_{B}  \frac{  |\phi(y)-\phi(x)|}{\mu(B_{x,y}) d(x,y)^s}  \,d\mu(x)
		\lesssim \frac{1}{s(1-s)}\frac{1}{(R-\rho)^s}.
		\]
		Inserting this into \eqref{eq2.3-13oct}, we obtain 
		{
			\[
			0\gtrsim   \int_B\int_B \frac{|\eta(x)-\eta(y)| }{\mu(B_{x,y}) d(x,y)^s}  \,d\mu(x)\,d\mu(y) - \frac{1}{s(1-s)}\frac{1}{(R-\rho)^s}\int_{B} (u-k)_+\, d\mu.
			\]
			Notice from the definition of subminimizers and $\mathcal K_{\psi,0}$-solution that $u\in \dot W^{s,1}(\Omega)$, and so $u$ is locally integrable on $\Omega$ by \eqref{eq2.1-11dec}. In particular, $\int_{B} (u-k)_+\, d\mu$ is finite, and hence by adding both sides of the above estimate  by the finite quantity $ \frac{1}{s(1-s)}
			\frac{1}{(R-\rho)^s}
			\int_{B} (u-k)_+\, d\mu$, 
			we obtain that
			\begin{align}
				\|(u-k)_+\phi\|_{\dot W^{s,1}(B(x_0,R))}
				\lesssim
				\frac{1}{s(1-s)}
				\frac{1}{(R-\rho)^s}
				\int_{B(x_0,R)} (u-k)_+\, d\mu.
				\label{eq2.1-19oct}
			\end{align}
			By the estimate  \eqref{eq2-3Lemma2}, it follows from $\phi\equiv 0$ on $X\setminus B(x_0,R)$ that
			\begin{align*}
				&\quad \|(u-k)_+\phi\|_{\dot W^{s,1}(X)}\\
				&= \|(u-k)_+\phi\|_{\dot W^{s,1}(B(x_0,R))} \\
				&\quad + 2 \int_{B(x_0,R)} \left( \int_{X\setminus B(x_0,R)} \frac{(u-k)_+(y)\phi(y)}{\mu(B_{x,y})d(x,y)^s}d\mu(x) \right)d\mu(y)\\
				&\overset{\eqref{eq2.1-19oct} \&  \eqref{eq2-3Lemma2}}{\lesssim} \frac{1}{s(1-s)}\frac{1}{(R-\rho)^s}\int_{B(x_0,R)} (u-k)_+\, d\mu\\
				&\quad + \frac{1}{s(1-s)}\int_{B(x_0,R)}  \frac{(u-k)_+(y)\phi(y)}{{( R- d(y,x_0))^s}}\, d\mu(y)\\
				& {\lesssim} \, \frac{1}{s(1-s)}\frac{1}{(R-\rho)^s}\int_{B(x_0,R)} (u-k)_+\, d\mu
			\end{align*}
			where the last estimate is obtained since $0\leq \phi\leq 1,$  $\phi\equiv  0$ on $ X\setminus B(x_0,r)$, and  $R-d(y,x_0)\geq R-r=\frac{R-\rho}{2}$ for $y\in B(x_0,r)$. The desired claim \eqref{eq3.2-15mar}  follows.
		}
	\end{proof}

\end{document}
